Esta investigación desde la perspectiva social, la reducción de accidentes de tránsito contribuye directamente a la disminución de mortalidad y lesiones graves. Desde el ámbito institucional, la disponibilidad de información analítica permite optimizar recursos y priorizar intervenciones. En cuanto a la ingeniería de software, este proyecto representa una aplicación práctica de arquitecturas escalables, bases de datos geoespaciales y analítica interactiva. Por lo anterior la implementación de la plataforma puede servir como base para desarrollos futuros en analítica predictiva y sistemas inteligentes de movilidad \cite{londono2024}.\\
La seguridad vial se ha consolidado como uno de los desafíos más relevantes para la gestión urbana contemporánea. El crecimiento acelerado de las ciudades, el aumento del parque automotor y la intensificación de los flujos de movilidad han incrementado la exposición al riesgo en los entornos urbanos. En este escenario, la accidentalidad de tránsito no solo representa un problema de salud pública, sino también un fenómeno social y económico que impacta la calidad de vida de las personas, genera pérdidas humanas irreparables y produce altos costos para los sistemas de salud, la infraestructura vial y la productividad de las ciudades \cite{lugo2025}.\\
En ciudades densamente pobladas como Medellín, la complejidad del sistema de movilidad exige herramientas tecnológicas que permitan comprender el comportamiento espacial y temporal de los siniestros viales. Aunque las instituciones suelen disponer de registros de accidentes, estos datos frecuentemente se encuentran dispersos en diferentes plataformas, formatos o repositorios, lo que dificulta su análisis integral y limita la capacidad de generar diagnósticos oportunos. Esta fragmentación de la información impide identificar patrones de riesgo, anticipar comportamientos críticos en determinadas zonas o evaluar con precisión el impacto de las medidas de seguridad implementadas \cite{castrillon2025}.\\
La investigación propone el desarrollo de un instrumento virtual basado en tecnologías geoespaciales y modelos de análisis de datos, capaz de integrar información relacionada con accidentes de tránsito en una plataforma digital estructurada. Este sistema permitirá almacenar registros de siniestros, clasificarlos según variables relevantes como ubicación geográfica, franja horaria, tipo de vehículo involucrado y nivel de gravedad y posteriormente procesarlos mediante herramientas de análisis estadístico \cite{calvo2025}.\\
A partir de esta estructura de datos, la plataforma incorporará mecanismos de visualización interactiva que faciliten la interpretación de la información mediante mapas georreferenciados, gráficos dinámicos y tableros de monitoreo. Estas herramientas permitirán identificar zonas críticas, tendencias temporales y comportamientos recurrentes de la accidentalidad en el territorio urbano. Asimismo, el sistema será diseñado bajo criterios de arquitectura de software que garanticen seguridad, escalabilidad y capacidad de actualización, permitiendo su adaptación a diferentes contextos urbanos o a futuras ampliaciones de la base de datos \cite{torres2022}.

\bigskip

A escala global, los siniestros viales continúan siendo una de las expresiones más persistentes de riesgo cotidiano en las ciudades: cada año mueren aproximadamente 1,19 millones de personas por accidentes de tránsito y entre 20 y 50 millones sufren lesiones no fatales, muchas con discapacidad, lo que revela un impacto humano sostenido que trasciende la estadística y afecta hogares, productividad y sistemas de salud \cite{ramirez2024}. Aunque existe evidencia de avances en reducción relativa frente a años previos, la magnitud del fenómeno mantiene la seguridad vial como un desafío de política pública y de gestión urbana basada en datos, especialmente porque las muertes se concentran de forma desproporcionada en países de ingresos bajos y medios y porque las lesiones afectan de manera sensible a población joven en edad productiva \cite{oms2024}. En este contexto, los sistemas que registran, integran y analizan información confiable se convierten en un componente estratégico: sin datos oportunos y comparables, las intervenciones se vuelven reactivas, fragmentadas y difíciles de evaluar.
En América Latina, el problema se expresa como una combinación de crecimiento del parque automotor, brechas de control y comportamientos de alto riesgo. En Argentina, se han descrito impactos relevantes en mortalidad y morbilidad asociada a siniestros, con afectación predominante en hombres y determinados grupos etarios, lo que refleja un patrón regional donde la carga se concentra en población económicamente activa \cite{aguilar2023}.\\

En México, se han impulsado estrategias preventivas desde el sector salud (por ejemplo, iniciativas de seguridad vial orientadas a reducir hospitalizaciones y muertes), evidenciando que la respuesta institucional tiende a fortalecerse cuando se cuenta con mecanismos de información y seguimiento \cite{ops2023}.\\

En Ecuador, reportes oficiales han atribuido una parte importante de los siniestros a imprudencia del conductor y han identificado tipos de vehículos frecuentemente involucrados, lo que indica que las políticas requieren segmentación por actor y causa para ser efectivas \cite{oviedo2025}.\\

En Perú, la concentración urbana también es determinante: Lima Metropolitana ha registrado proporciones altas de accidentes respecto al total nacional, lo que refuerza que la ciudad es el escenario crítico donde se cruzan volumen de tráfico, exposición y riesgo 
En Colombia, la siniestralidad vial sigue siendo un problema estructural con víctimas fatales y lesionados año a año, y con presiones significativas sobre el sistema sanitario y la seguridad urbana. Para 2023 se reportaron preliminarmente 8.405 personas fallecidas por siniestros viales a nivel nacional, lo que dimensiona la urgencia de fortalecer acciones preventivas y de control con enfoque territorial \cite{betancourt2021}. Además, se reconoce que la carga puede concentrarse en departamentos y ciudades con alta movilidad; y, en el debate público, Medellín aparece de manera recurrente entre capitales donde se reflejan estas dinámicas \cite{valencia2025}.\\

El problema central, por tanto, no es únicamente la ocurrencia de siniestros, sino la forma en que la información se captura, integra y convierte en decisiones. En muchos contextos, los datos permanecen dispersos en reportes no interoperables, bases heterogéneas o consolidaciones tardías; esto limita la capacidad de detectar patrones en tiempo real, anticipar picos de riesgo, priorizar controles y evaluar si una intervención redujo realmente los eventos en un punto específico. Para Medellín, donde existen corredores con concentración histórica y dinámicas urbanas cambiantes, se vuelve crítico contar con un sistema digital que permita registrar siniestros de forma estructurada, georreferenciarlos, visualizar mapas de calor, producir estadísticas y habilitar tableros (dashboards) que apoyen decisiones de mitigación basadas en evidencia. En consecuencia, se justifica un proyecto de ingeniería de software orientado a diseñar y desarrollar una plataforma digital actualizable, con analítica geoespacial y reportes automatizados, capaz de transformar el dato operativo en inteligencia pública para reducir el impacto humano y social de la siniestralidad vial \cite{aguilar2023}.
La problemática de la accidentalidad vial no solo debe analizarse desde la ocurrencia de eventos en las vías, sino también desde sus múltiples impactos sociales, económicos e institucionales. Cada accidente de tránsito implica una serie de consecuencias que trascienden el momento del siniestro y se proyectan hacia distintos ámbitos de la sociedad. Desde una perspectiva social, los accidentes generan afectaciones profundas en las familias de las víctimas, particularmente cuando se trata de personas en edad laboral activa. La pérdida de un miembro del hogar o la presencia de lesiones incapacitantes puede alterar significativamente las dinámicas familiares, reducir los ingresos económicos y generar procesos prolongados de rehabilitación física y emocional \cite{torres2022}.\\
Por ende, la accidentalidad vial representa una presión constante para los sistemas de salud. Los centros hospitalarios reciben diariamente pacientes con traumas derivados de accidentes de tránsito, lo que implica la utilización de recursos médicos, quirúrgicos y de rehabilitación que podrían destinarse a otras necesidades sanitarias. En muchos casos, las lesiones derivadas de los accidentes requieren tratamientos prolongados, hospitalizaciones extensas o procesos de recuperación complejos que generan costos elevados para los sistemas de aseguramiento en salud y para las instituciones hospitalarias. De esta manera, la siniestralidad vial no solo afecta a las víctimas directas, sino que también genera impactos estructurales en la sostenibilidad del sistema sanitario \cite{follana2024}.\\
Desde una perspectiva económica, los accidentes de tránsito producen pérdidas significativas relacionadas con daños materiales, interrupción de actividades productivas y costos asociados a la atención médica y a los procesos legales. Las empresas, por ejemplo, pueden verse afectadas cuando sus trabajadores sufren lesiones que les impiden continuar con sus labores, lo que repercute en la productividad y en la continuidad de las operaciones. De igual manera, los daños a la infraestructura vial y a los vehículos implican gastos de reparación y reposición que afectan tanto a los ciudadanos como a las instituciones encargadas del mantenimiento del espacio público \cite{chavez2024}.\\
En el ámbito urbano, los accidentes de tránsito también generan impactos en la movilidad y en el funcionamiento de la ciudad. Un siniestro vial puede provocar congestiones prolongadas, interrupciones en el flujo vehicular y alteraciones en los tiempos de desplazamiento de miles de personas. En ciudades con alta densidad poblacional y sistemas de transporte complejos, estas interrupciones pueden tener efectos en cadena que afectan el funcionamiento de la red de movilidad urbana. Por esta razón, comprender la distribución espacial y temporal de los accidentes se convierte en un elemento fundamental para mejorar la planificación del tránsito y optimizar la gestión de la movilidad \cite{betancourt2021}.\\
Otro aspecto relevante se relaciona con la dimensión institucional de la seguridad vial. Las autoridades encargadas de la gestión del tránsito requieren información confiable y actualizada para diseñar estrategias de prevención y control. Sin embargo, en muchos contextos urbanos la información disponible sobre accidentes de tránsito se encuentra fragmentada entre diferentes entidades, como organismos de tránsito, cuerpos de policía, instituciones de salud y aseguradoras. Esta dispersión de datos dificulta la construcción de diagnósticos integrales y limita la capacidad de generar análisis que permitan comprender con precisión las causas y patrones de la accidentalidad \cite{ortiz2024}.
La ausencia de sistemas integrados de información también limita el desarrollo de estrategias de prevención basadas en evidencia. Cuando los datos sobre accidentes no se encuentran estructurados o georreferenciados adecuadamente, resulta más complejo identificar tendencias, analizar comportamientos recurrentes o evaluar el impacto de las medidas implementadas para mejorar la seguridad vial. En consecuencia, muchas decisiones relacionadas con la gestión del tránsito se basan en información parcial o en análisis retrospectivos que no siempre reflejan las dinámicas actuales de la movilidad urbana \cite{oviedo2025}.\\
En este contexto, el uso de herramientas tecnológicas y sistemas de análisis de datos se presenta como una alternativa clave para fortalecer la gestión del riesgo vial. Las tecnologías de información permiten integrar grandes volúmenes de datos provenientes de diferentes fuentes, procesarlos mediante modelos analíticos y representarlos de manera visual para facilitar su interpretación. Particularmente, los sistemas de información geoespacial han demostrado ser herramientas eficaces para analizar fenómenos urbanos complejos, ya que permiten representar la información en mapas interactivos que evidencian la distribución territorial de los eventos \cite{carrillo2024}.\\
La incorporación de análisis geoespacial en el estudio de la accidentalidad vial permite identificar zonas críticas donde los accidentes se concentran con mayor frecuencia, así como analizar las características específicas de estos eventos en términos de tipo de vehículo, condiciones de la vía, franjas horarias y gravedad de las lesiones. Este tipo de análisis facilita la comprensión de los factores que influyen en la ocurrencia de los siniestros y permite orientar intervenciones más precisas en los lugares donde el riesgo es mayor \cite{monar2025}.\\
En ciudades con grandes concentraciones urbanas, como Medellín, la implementación de herramientas tecnológicas para el análisis de la accidentalidad vial adquiere especial relevancia. La complejidad de su sistema de movilidad, la diversidad de actores viales y la intensidad del tráfico generan un escenario donde la gestión de la seguridad vial requiere enfoques cada vez más sofisticados. En este sentido, la disponibilidad de plataformas digitales capaces de integrar datos georreferenciados, realizar análisis estadísticos y generar visualizaciones dinámicas puede contribuir significativamente a mejorar la comprensión del fenómeno de la accidentalidad  \cite{parraga2024}.\\
Además, el desarrollo de sistemas digitales orientados al análisis de accidentes de tránsito se alinea con las tendencias actuales de transformación digital en la gestión urbana. Las ciudades contemporáneas buscan incorporar tecnologías de información para mejorar la eficiencia de los servicios públicos, optimizar la planificación del territorio y fortalecer la toma de decisiones basada en datos. En este marco, el uso de plataformas analíticas y sistemas de monitoreo digital permite avanzar hacia modelos de gestión más inteligentes y orientados a la prevención de riesgos \cite{changotasig2023}.\\
Por lo tanto, la necesidad de contar con herramientas tecnológicas que permitan analizar y visualizar la información sobre accidentes de tránsito no responde únicamente a un interés académico o investigativo, sino a una demanda concreta de las ciudades que buscan mejorar sus sistemas de movilidad y reducir el impacto de la siniestralidad vial. El desarrollo de un instrumento virtual geoespacial orientado a visualizar y analizar índices de accidentalidad representa una oportunidad para fortalecer la gestión del riesgo vial en contextos urbanos complejos  \cite{parraga2024}.\\
En consecuencia, el diseño e implementación de un sistema de mitigación de accidentes de tránsito para ciudades con grandes urbes, como el caso de Medellín, puede contribuir a transformar los datos sobre siniestros viales en información estratégica para la toma de decisiones. La integración de herramientas de análisis estadístico, visualización interactiva y georreferenciación permitiría identificar patrones de riesgo, monitorear tendencias y apoyar la formulación de estrategias de prevención orientadas a reducir la ocurrencia de accidentes y sus consecuencias en la sociedad \cite{figueroa2022}.\\
